\documentclass[11pt]{amsart}
\usepackage[margin=1in]{geometry}
\usepackage{amsmath,amssymb,amsthm,booktabs,tikz,graphicx}
\usetikzlibrary{positioning}
\usepackage[pdfusetitle]{hyperref}
\pdfstringdefDisableCommands{\def\\{ }\def\times{x}}
\usepackage[capitalize]{cleveref}

\newtheorem{theorem}{Theorem}[section]
\newtheorem{lemma}[theorem]{Lemma}
\newtheorem{proposition}[theorem]{Proposition}
\newtheorem{corollary}[theorem]{Corollary}
\newtheorem{conjecture}[theorem]{Conjecture}
\newtheorem{problem}[theorem]{Problem}
\crefname{problem}{Problem}{Problems}
\theoremstyle{definition}
\newtheorem{definition}[theorem]{Definition}
\newtheorem{remark}[theorem]{Remark}
\newtheorem{heuristic}[theorem]{Heuristic}

\newcommand{\Ham}{\mathrm{H}}
\newcommand{\Zhat}{\widehat{\mathbb{Z}}}
\newcommand{\KL}{\mathrm{KL}}

\title[Hamiltonian paths in square-sum graphs]{Hamiltonian paths in square-sum graphs:\\
an exponential lower bound, a limit object,\\ and a conjectural asymptotic}
\author{Victor S. Miller}
\date{Draft, \today}

\begin{document}

\begin{abstract}
Let $G(n)$ be the graph on $\{1,\dots,n\}$ in which $i\sim j$ iff $i+j$ is a perfect
square, and let $\Ham(n)$ be the number of its (undirected) Hamiltonian paths.
Gerbicz showed that $\Ham(n)>0$ for all $n\ge 25$. We isolate a class of local
\emph{parity-preserving moves} which are multiplied by $25$ under Gerbicz's
$\times 25$ lifting step, whereas path counts are not, and deduce the exponential bound
$\Ham(N)\ge 2^{(4N+2)/91}$ along the sequence $N=(91\cdot 25^m-1)/2$,
from a computer-verified seed in $G(45)$. We explain why $25$ is the smallest
square factor for which the lifting scheme closes, and discuss what is needed to
reach all large $N$. We then turn to limits. The rescaled graphs do not converge in cut
norm to the obvious kernel $1/(2\sqrt{x+y})$: the arithmetic of squares
survives in the limit, which naturally lives on $[0,1]\times\Zhat$ as a singular
measure. A first-moment computation in this limit predicts
$\log \Ham(n) = \log n! - \tfrac n2\log n + Cn + o(n)$ with
$C=-0.632538\ldots$ given by a Schr\"odinger (entropic transport) problem. This is
consistent with exact counts for $15\le n\le 74$ to within $O(1)$, and lies
just below the rigorous Bregman upper bound, whose constant is $B=-0.614954\ldots$.
The results of \cref{sec:lift,sec:moves} are theorems
(with computer-verified certificates); those of \cref{sec:limits,sec:heuristic,sec:superexp}
are partly heuristic, and we mark clearly which statements are proved.
\end{abstract}

\maketitle

\section{Introduction}

For $n\ge1$ let $G(n)$ have vertex set $[n]=\{1,\dots,n\}$, with $\{i,j\}$ an edge iff
$i+j$ is a perfect square. Hamiltonian paths of $G(n)$ are the \emph{square chains}
of recreational mathematics. They exist exactly for $n\in\{15,16,17,23\}$ and all $n\ge25$
(OEIS \cite{OEIS-A090461}). The existence for all $n\ge 25$, and of Hamiltonian cycles
for all $n\ge32$, is due to Gerbicz \cite{Gerbicz}, by a combination of
construction and computation; see also \cite[Appendix]{BerlekampGuy}.
The number $\Ham(n)$ of undirected Hamiltonian paths is OEIS A071983
\cite{OEIS-A071983}. Exact values are known for $n\le 75$; they were computed by a frontier dynamic program along an exact pathwidth ordering
(cf.\ \cite{Pettersson,Coudert}); we use the values in \texttt{data/b071983.txt}.

Since $G(n)$ is an induced subgraph of $G(n+1)$, the graphs form a nested
sequence with direct limit $G(\infty)$ on $\mathbb{N}$. Every vertex $i\le n$ has
$\sqrt{n+i}-\sqrt{i}+O(1)=\Theta(\sqrt n)$ neighbours, so $G(n)$ has edge density
$\asymp n^{-1/2}$. This is the ``intermediate'' regime, too sparse for graphons and too
dense for local (Benjamini--Schramm) limits.

\subsection*{Results}
\begin{enumerate}
\item (\cref{thm:main}) $\Ham(N_m)\ge 2^{2\cdot 25^m}$ for $N_m=(91\cdot25^m-1)/2$,
i.e.\ $\Ham(N)\ge 2^{(4N+2)/91}$ along this sequence. The proof combines Gerbicz's
lifting step with a notion of \emph{parity-preserving move} (\cref{def:move}),
which is the invariant that survives lifting with full multiplicity.
\item (\cref{prop:t3}) Within the lifting scheme, the factor $9$ is impossible
and $25$ is the smallest square that works. There is also a heuristic reason
(\cref{rem:threshold}).
\item (\cref{prop:mod4}) The rescaled graphs $\sqrt n\,G(n)$ do not converge in cut norm to
$\kappa(x,y)=1/(2\sqrt{x+y})$. We propose a limit on $[0,1]\times\Zhat$
(\cref{conj:limit}).
\item (\cref{conj:asymp}) $\log\Ham(n)=\log n!-\frac n2\log n+Cn+o(n)$ with
$C\approx-0.632538$. The prediction is compared with data in \cref{tab:resid}. Rigorously,
$\log\Ham(n)\le\frac n2\log n+(B-1)n+o(n)$ with
$B=\frac{\sqrt2}2-\frac32\log(1+\sqrt2)\approx-0.61495$, and $C<B$ (\cref{prop:bregman}).
\item (\cref{sec:deterministic}) The count of cycle covers transfers to $G(n)$
deterministically: by Gurvits's Bethe bound, any doubly stochastic matrix supported on
$E(G(n))$ gives a lower bound for $\log\operatorname{per}A_{G(n)}$, and the Sinkhorn scaling
of $A_{G(n)}$ reproduces $C-1$ numerically up to $n=64000$. \Cref{app:limits} explains graph
limits, permutons and the Schr\"odinger problem in detail.
\item (\cref{sec:superexp}) If \cref{conj:asymp} is true, $\Ham(n)$ is superexponential.
We outline the permanent-plus-rotation method of
\cite{Krivelevich,GlebovKrivelevich} as a possible route to bounds of that order.
\end{enumerate}

\section{Gerbicz's lifting step}\label{sec:lift}

Throughout, $t$ is odd, $T=t^2$, and $h=(T-1)/2$. The basic observation is that
$T(a+b)=(tk)^2$ whenever $a+b=k^2$. We give a self-contained account of the
$\times25$ step that underlies \cite{Gerbicz}. (The glue arrangement in
\cref{tab:glue} was found independently by computer search and may differ from
Gerbicz's.)

\begin{lemma}[Lifted segments]\label{lem:segments}
Let $b=(b_1,\dots,b_n)$ be a Hamiltonian path of $G(n)$. For $c\in[-h,h]$ put
\[
  T_c(b)=\bigl(T b_p+(-1)^{p+1}c\bigr)_{p=1}^n .
\]
Then each $T_c(b)$ is a path in $G(Tn+h)$. The vertex sets of the $2h+1$ paths
$T_c(b)$ partition $[h+1,\,Tn+h]$.
\end{lemma}

\begin{proof}
Consecutive terms have opposite signs on $c$, so their sum is
$T(b_p+b_{p+1})$, which is a square. Each integer in $[h+1,Tn+h]$ is uniquely of the
form $Tb+c$ with $1\le b\le n$ and $|c|\le h$. For $c\ne0$, the number $Tb+c$ lies in
$T_c(b)$ if $b$ sits at an odd position of $b$, and in $T_{-c}(b)$ if it sits at an even
position, and in no other $T_{c'}(b)$. The numbers $Tb$ all lie in $T_0(b)$.
\end{proof}

\begin{remark}\label{rem:pairs}
\Cref{lem:segments} remains true if $T_c$ and $T_{-c}$ are built from two different
base paths $b,b'$, provided $b$ and $b'$ have the same set of vertices at even
positions. Only the $h+1$ classes $\{0\},\{\pm1\},\dots,\{\pm h\}$ can choose
their base paths independently. Counting paths therefore gives only
$\log\Ham(Tn+h)\gtrsim (h+1)\log\Ham'(n)$, where $\Ham'$ counts base paths with the
required endpoints. This yields $\log\Ham(N)\gtrsim N^{\alpha}$ with
$\alpha=\log((T+1)/2)/\log T<1$ (for $t=5$, $\alpha=\log13/\log25\approx0.797$).
\Cref{sec:moves} removes this loss.
\end{remark}

\subsection{Gluing}
Suppose now that $n$ is odd and $b$ runs from $b_1=1$ to $b_n=3$. Then $T_c(b)$ runs from
$T+c$ to $3T+c$. These endpoints are independent of $n$. A \emph{glue
arrangement} is a Hamiltonian path from $1$ to $3$ in the auxiliary graph on
$\{1,\dots,h\}\cup\{T+c,3T+c:|c|\le h\}$ (edges: square sums) which traverses each
pair $(T+c,3T+c)$ consecutively. Replacing each such pair by $T_c(b)$, or by its reverse,
gives:

\begin{proposition}\label{prop:lift}
If a glue arrangement exists for $t$, then for every odd $n$ and every Hamiltonian
path of $G(n)$ from $1$ to $3$, there is a Hamiltonian path of $G(Tn+h)$ from $1$ to $3$
that contains each $T_c(b)$ as a contiguous block. Moreover $Tn+h$ is odd if $h$ is even.
\end{proposition}

For $t=5$ we found the glue arrangement of \cref{tab:glue}. Iterating from a seed $n_0$
gives Hamiltonian paths from $1$ to $3$ in $G(N_m)$ with
$N_{m+1}=25N_m+12$, i.e.\ $N_m=(n_0+\frac12)25^m-\frac12$. With
$n_0=35$ this is Gerbicz's sequence $(71\cdot25^m-1)/2$ \cite{OEIS-A090461}.

\begin{table}[h]
\small
\[
\begin{array}{l}
1\ [15\ 65]\ [35\ 85]\ [36\ 86]\ [14\ 64]\ [80\ 30]\ [70\ 20]\ 5\ 11\ [25\ 75]\ 6\ [19\ 69]\ 12\ [24\ 74]\ 7\\
{}[29\ 79]\ 2\ [34\ 84]\ [37\ 87]\ [82\ 32]\ 4\ [21\ 71]\ 10\ [26\ 76]\ [68\ 18]\ [31\ 81]\ [63\ 13]\ [23\ 73]\ 8\\
{}[28\ 78]\ [66\ 16]\ 9\ [72\ 22]\ [27\ 77]\ [67\ 17]\ [83\ 33]\ 3
\end{array}
\]
\caption{A glue arrangement for $t=5$. A bracket $[x\ y]$ denotes the segment
$T_c$ with endpoints $\{x,y\}=\{25+c,75+c\}$, traversed from $x$ to $y$ (so
$[80\ 30]$ is $T_5$ reversed). Unbracketed entries are glue vertices. Each of the
$36$ junctions, i.e.\ each pair of consecutive entries not inside one bracket, sums to a square:
$1+15=16$, $65+35=100$, $85+36=121$, \dots, $33+3=36$.}
\label{tab:glue}
\end{table}

\begin{figure}[h]
\centering
\resizebox{\linewidth}{!}{\input{../figs/gluing.tex}}
\caption{The start of the glued path in $G(1137)$, lifted from the seed
$b=(1,35,29,20,\dots,33,3)$ in $G(45)$. Inside a segment, consecutive sums are
$25\times$ the corresponding sums in $b$. The junction between $T_{-10}$ and $T_{10}$ is
the ordinary square $65+35=100$.}
\label{fig:gluing}
\end{figure}

\subsection{Why $25$?}
The multiplier must be a square $m=t^2$: $mk^2$ must be a square for all $k$. The
alternating offsets $\pm c$ then use one representative of each residue class mod $T$.
For even $t$ the class of $T/2$ is self-paired, and covering each $T(k+\frac12)$
exactly once would force all vertices of $b$ to occupy positions of the same parity,
which is impossible. So $t$ is odd. For the step to iterate, the output endpoints
must again be glue numbers. So the base path runs between some $u,w\le h$, and
the glue problem is a fixed finite problem for each $(t,u,w)$ and each parity of $n$.

\begin{proposition}[computer-verified]\label{prop:t3}
For $t=3$ there is no glue arrangement for any $u\ne w$ in $\{1,2,3,4\}$ and either
parity of $n$ (exhaustive search, 24 cases). For $t=5$ arrangements exist for
$(u,w)=(1,3)$ and $(1,2)$ with both parities. For $t=7$ one was found for $(1,2)$
with $n$ even.
\end{proposition}

\begin{remark}\label{rem:threshold}
Heuristically, the endpoints of the lifted segments lie in windows of width $\approx T$
around $Tu$ and $Tw$. Most junctions join two such windows, with sums near $T(u+w)$, where
consecutive squares are $\approx 2t\sqrt{u+w}$ apart. Enough square-sum pairs are
available only if $T\gtrsim 2t\sqrt{u+w}$, i.e.\ $t>2\sqrt{u+w}$. With $(u,w)=(1,3)$ this
is $t>4$, and for $t=3$ it already fails at $u+w=3$. \Cref{prop:t3} concerns only this
scheme (glue set $\{1,\dots,h\}$, base endpoints $\le h$). Other schemes with
multiplier $9$ are not excluded.
\end{remark}

\section{Moves that survive lifting}\label{sec:moves}

\begin{definition}\label{def:move}
Let $P=(v_1,\dots,v_N)$ be a Hamiltonian path. A \emph{move} of $P$ is a window
$[i,j]$ with $2\le i\le j\le N-1$, together with an ordering
$(w_i,\dots,w_j)\ne(v_i,\dots,v_j)$ of $\{v_i,\dots,v_j\}$ such that
$v_{i-1},w_i,\dots,w_j,v_{j+1}$ is a path, and such that each vertex keeps the
parity of its position: if $w_p=v_q$ then $p\equiv q\pmod 2$. A family of moves is
\emph{separated} if the windows are pairwise disjoint and any two are at
distance at least $2$ (some position lies strictly between them).
\end{definition}

Odd-length segment reversals with the two chords present are moves. So are
transpositions of two adjacent blocks of even length. A pair of P\'osa rotations
\cite{Posa} at the same end that restores the endpoint is a fixed-endpoint rerouting;
it is a move if it preserves parities.

\begin{lemma}\label{lem:2k}
If $P$ has $k$ separated moves then $G$ has at least $2^k$ Hamiltonian paths with the
same endpoints as $P$.
\end{lemma}
\begin{proof}
Separation means the anchors $v_{i-1},v_{j+1}$ of one move are not inside another
window. So any subset $S$ of the moves can be applied simultaneously, giving a
Hamiltonian path $P_S$. Different subsets give different sequences, since the content
of each window records whether its move was applied. All $P_S$ start at $v_1$ and end at
$v_N$, so they are also distinct as undirected paths.
\end{proof}

\begin{lemma}[Lifting moves]\label{lem:liftmoves}
In the situation of \cref{prop:lift}, every move of $b$ induces a move of the lifted path
$P'$ inside each of the $2h+1$ blocks $T_c(b)$. If the moves of $b$ are separated,
so are the resulting $(2h+1)k$ moves of $P'$.
\end{lemma}
\begin{proof}
Let $(w_i,\dots,w_j)$ reorder $(b_i,\dots,b_j)$ with parities preserved. In $T_c(b)$,
replace $Tb_p+(-1)^{p+1}c$ ($i\le p\le j$) by $Tw_p+(-1)^{p+1}c$. Since $w_p$
and the vertex it replaces have positions of the same parity, the new entries form
the same set. Consecutive sums, including the two at the anchors, are
$T\cdot(\text{square})$. In $P'$ the block $T_c(b)$ occupies consecutive positions,
traversed forwards or backwards. Since $n$ is odd, reversal $p\mapsto n+1-p$
preserves parity. Hence positions within the block are shifted by a constant
(possibly after reversal), parities are uniformly preserved or uniformly flipped,
and the induced reordering preserves parities in $P'$. The windows are interior to
$b$, so the block endpoints and the glue are untouched. Separation is preserved
within a block. Windows in different blocks are separated because each window is
interior to its block, so the block endpoints lie between them.
\end{proof}

The point is that a parity-preserving move does not change the even-position set,
so $T_c$ and $T_{-c}$ may apply it independently. This removes the pairing loss of
\cref{rem:pairs}.

\begin{theorem}\label{thm:main}
Let $n_0$ be odd and suppose $G(n_0)$ has a Hamiltonian path from $1$ to $3$ with $k$
separated moves. Let $N_m=(n_0+\frac12)25^m-\frac12$. Then $G(N_m)$ has a
Hamiltonian path from $1$ to $3$ with $k\,25^m$ separated moves, and
\[
\Ham(N_m)\ \ge\ 2^{k\,25^m}\ =\ 2^{k(2N_m+1)/(2n_0+1)} .
\]
\end{theorem}
\begin{proof}
Induction on $m$ using \cref{prop:lift} with \cref{tab:glue}, and
\cref{lem:liftmoves}; then \cref{lem:2k}.
\end{proof}

\begin{corollary}[computer-verified seeds]\label{cor:seeds}
\begin{enumerate}
\item $n_0=45$, $k=2$:\quad $\Ham(N)\ge 2^{(4N+2)/91}$ for $N=(91\cdot25^m-1)/2$
$=45,1137,28437,\dots$
\item $n_0=35$, $k=1$:\quad $\Ham(N)\ge 2^{(2N+1)/71}$ for $N=(71\cdot25^m-1)/2$
$=35,887,22187,\dots$ (Gerbicz's sequence).
\item $n_0=43$, $k=1$:\quad $\Ham(N)\ge 2^{(2N+1)/87}$ for $N=(87\cdot25^m-1)/2$.
\end{enumerate}
\end{corollary}

The seed for (1) is
\[
\begin{aligned}
&1\ 35\ 29\ 20\ 44\ 37\ 27\ 22\ 14\ 11\ 5\ 31\ 18\ 7\ 42\ 39\ 25\ 24\ 12\ 13\ 36\ 28\ 8\ 17\ 32\ 4\ 45\ 19\ 30\\
&6\ 10\ 26\ 38\ 43\ 21\ 15\ 34\ 2\ 23\ 41\ 40\ 9\ 16\ 33\ 3,
\end{aligned}
\]
with the two moves
\[
\begin{aligned}
&1\,35\,|\,29\,20\,44\,37\,27\,22\,|\,14\,11\,5\,31\,18\,7\,|\,42\\
&\quad\longmapsto\ 1\,35\,|\,14\,11\,5\,31\,18\,7\,|\,29\,20\,44\,37\,27\,22\,|\,42,\\[2pt]
&28\,|\,8\,17\,32\,4\,45\,19\,30\,6\,10\,26\,38\,43\,21\,15\,34\,2\,23\,|\,41\\
&\quad\longmapsto\ 28\,|\,21\,43\,38\,26\,23\,2\,34\,15\,10\,6\,30\,19\,45\,4\,32\,17\,8\,|\,41 .
\end{aligned}
\]
(The first is a transposition of two blocks of length $6$, with chords
$35+14=49$, $7+29=36$, $22+42=64$. The second reverses the block of length $17$ between $28$ and $41$ and then
exchanges the two blocks $23\,2\,34\,15$ and $21\,43\,38\,26$ inside it; the new
adjacencies are $28+21=49$, $26+23=49$, $15+10=25$ and $8+41=49$.) The program
\texttt{hamiltonian.bounds.certificate} checks the seeds and builds the lifts
explicitly for $m=1,2$ ($N=1137$ and $28437$), verifying the $50$ and $1250$
separated moves. Exhaustive search over all Hamiltonian paths from $1$ to $3$ gives
$k\le1$ for all odd $n_0\le 41$ (with equality for $n_0=35,39,41$). Random sampling
(windows of length $\le24$) gives $k=2$ at $n_0=45,55,63,65,69$, and $n_0=45$ has the largest ratio
$k/(n_0+\frac12)$ found so far.

\begin{remark}
A seed with $k\ge1$ is not actually needed. Any level at which the glue or the
``cross'' edges (sums that are not $25\times$ a square) create one new move
gives the same conclusion, because every later level multiplies it by $25$.
\end{remark}

\subsection{All large $N$}
\Cref{thm:main} gives exponential growth only along sparse sequences. Gerbicz's
general construction \cite{Gerbicz} reaches every $N\ge25$ using further offsets and
residues. Write $k(N)$ for the largest number of separated moves in a Hamiltonian path of
$G(N)$ with the endpoints that construction requires.

\begin{proposition}[conditional]\label{prop:allN}
Suppose that for all large $N$ there are $n$ and $t$ (from a finite set of odd
$t\ge5$) with $N\le t^2 n+O(t^2)$, and a Hamiltonian path of $G(N)$ containing, as
contiguous blocks, $t^2$ lifted segments $T_c(b^{(c)})$ whose base paths $b^{(c)}$
are Hamiltonian paths of $G(n_c)$ with $n_c\in\{n-1,n,n+1\}$ and the required endpoints.
Then $\liminf k(N)/N>0$, and $\Ham(N)\ge 2^{\delta N}$ for some $\delta>0$ and all
large $N$, provided the finitely many base cases in a range $[n_*,t^2n_*]$ have
$k(n)\ge\delta_0 n$.
\end{proposition}
\begin{proof}[Sketch]
By \cref{lem:liftmoves}, $k(N)\ge t^2\min_c k(n_c)$. With $\rho(N)=k(N)/N$ this gives
$\rho(N)\ge\rho(n')(1-O(t^2/N))$ for some $n'\approx N/t^2$. Along any descending
chain the $N$'s grow geometrically, so $\prod(1-O(t^2/N_j))$ converges to a
positive limit.
\end{proof}
We have not checked that Gerbicz's general step has this ``every class is a lift''
form. This is the main gap between \cref{thm:main} and an exponential bound for all $N$.

\section{Limit objects}\label{sec:limits}

\begin{proposition}\label{prop:edges}
$e(G(n))=\frac23(\sqrt2-1)\,n^{3/2}+O(n)$.
\end{proposition}
\begin{proof}
$e(G(n))=\sum_{k}\#\{i<j\le n:\ i+j=k^2\}$, and the summand is $k^2/2+O(1)$ for
$k^2\le n$ and $n-k^2/2+O(1)$ for $n<k^2\le 2n$. Summing gives
$\frac{n^{3/2}}6+(\sqrt2-1)n^{3/2}-\frac{2\sqrt2-1}6n^{3/2}+O(n)$.
\end{proof}
Equivalently, $e(G(n))\sim\frac{n^{2}}{2}\cdot n^{-1/2}\iint_{[0,1]^2}\kappa$, with
\[
\kappa(x,y)=\frac{1}{2\sqrt{x+y}},
\]
the local density of squares near $n(x+y)$. Note $\kappa\in L^p([0,1]^2)$ exactly for $p<4$.
The natural candidate limit, in the sense of the $L^p$ graphons of
Borgs--Chayes--Cohn--Zhao \cite{BCCZ1,BCCZ2}, is therefore $\kappa$ with
normalisation $\rho_n=n^{-1/2}$. It is the wrong one:

\begin{proposition}\label{prop:mod4}
Let $S_n=4\mathbb{Z}\cap[1,n]$. Then $e(G(n)[S_n])\sim\frac18 e(G(n))$, whereas
$\kappa$ predicts $\frac1{16}e(G(n))$. Consequently, for the step graphon $W_n$ of
$\sqrt n\,A_{G(n)}$ (natural labelling), $\|W_n-\kappa\|_\square\not\to0$.
\end{proposition}
\begin{proof}
$a\mapsto 4a$ is an isomorphism $G(\lfloor n/4\rfloor)\to G(n)[S_n]$, because
$4a+4b$ is a square iff $a+b$ is. By \cref{prop:edges},
$e(G(n)[S_n])\sim\frac23(\sqrt2-1)(n/4)^{3/2}=\frac18e(G(n))$. Because $S_n$ is
equidistributed in $[0,1]$ after scaling, $\kappa$ predicts the fraction
$(1/4)^2=1/16$. The sets $S_n/n$ give a cut-norm discrepancy bounded below.
\end{proof}
(\texttt{hamiltonian.bounds.edges} confirms the ratio $2.0000$ at $n=10^8$.)

\begin{remark}[failure of upper regularity; heuristic]
Partition $[n]$ by residues modulo $Q$. The averaged density on block $(r,s)$ is
asymptotically $\varrho_Q(r+s)$ times the $\kappa$-prediction, where
$\varrho_Q(a)=\#\{x\bmod Q:x^2\equiv a\}$. For odd primes $\ell\mid Q$, half the
residues have $\varrho_\ell=2$, so the $L^p$ norm of the block-averaged graphon is about
$\prod_{\ell\mid Q}\bigl(2^{p-1}+O(1/\ell)\bigr)$, which is unbounded as $Q$ runs
through primorials. So the sequence is not $L^p$ upper regular in the sense of
\cite{BCCZ1} for any $p>1$, and that theory does not apply directly.
\end{remark}

\subsection{A profinite limit}
Map vertex $i$ of $G(n)$ to $(i/n,\,i)\in\Omega=[0,1]\times\Zhat$, with measure
$\lambda\otimes\mu_{\Zhat}$ (Lebesgue times Haar). Let $\sigma$ be the law of $k^2$
for Haar-random $k\in\Zhat$, and let $\nu$ be the law of $(\alpha,\beta)$ with $\alpha$
Haar and $\beta=s-\alpha$, $s\sim\sigma$ independent. Since squares form a Haar-null
set, $\nu$ is singular with respect to $\mu_{\Zhat}^{\otimes2}$. Define the symmetric
measure
\[
\Lambda=\kappa(x,y)\,dx\,dy\ \otimes\ \nu(d\alpha,d\beta)\qquad\text{on }\Omega^2 .
\]
By the equidistribution of squares in residue classes, for intervals $I,J$ and residue
classes $A,B$ modulo $Q$,
\[
\frac{\#\{(i,j)\in E(G(n)):\ i/n\in I,\ j/n\in J,\ i\in A,\ j\in B\}}{n^{3/2}}
\ \longrightarrow\ \Lambda\bigl((I\times A)\times(J\times B)\bigr).
\]
So $\Lambda$ captures all ``finite-level'' statistics. Measure-valued limits of
this kind (s-graphons) are the subject of Kunszenti-Kov\'acs, Lov\'asz and Szegedy
\cite{KLS}. The more general framework of action convergence and graphops is due to
Backhausz and Szegedy \cite{BackhauszSzegedy}; it covers graphons, $L^p$ graphons and
graphings, and is designed for intermediate densities.

\begin{conjecture}\label{conj:limit}
$n^{-1/2}A_{G(n)}$ action-converges, in the sense of \cite{BackhauszSzegedy}, to the
graphop on $(\Omega,\lambda\otimes\mu_{\Zhat})$ with kernel measure $\Lambda$.
\end{conjecture}

The archimedean part is scale-covariant: $\kappa(sx,sy)=s^{-1/2}\kappa(x,y)$.
This is reminiscent of graphex models for sparse graphs \cite{VeitchRoy,BCCH}, but
$G(\infty)$ is not exchangeable.

\section{A conjectural asymptotic for $\Ham(n)$}\label{sec:heuristic}

\begin{heuristic}\label{heu:first}
Consider the random graph $R(n)$ with independent edges of probability
$p_{ij}=\kappa(i/n,j/n)/\sqrt n$. A Hamiltonian cycle is a cyclic permutation $\pi$,
weighted by $\prod_i p_{i\pi(i)}$. To exponential order, restricting from all
permutations to cyclic ones and from cycles to paths costs only polynomial factors.
So
\[
\mathbb{E}\,\Ham(R(n))\ =\ \exp\Bigl(\log n!-\tfrac n2\log n+Cn+o(n)\Bigr),
\qquad
C=\lim_{n\to\infty}\tfrac1n\log\mathbb{E}_{\pi}\prod_{i=1}^n\kappa\bigl(\tfrac in,\tfrac{\pi(i)}n\bigr),
\]
with $\pi$ uniform in $S_n$, and by the large-deviation principle for permutons
(\cite[Theorem~1]{KKRW}, restated from Trashorras),
\[
C=\sup_{\mu}\Bigl\{\int\log\kappa\,d\mu-\KL(\mu\,\|\,\lambda\otimes\lambda)\Bigr\},
\]
with the supremum over couplings $\mu$ of two uniform marginals (permutons).
The unboundedness of $\log\kappa$ near $(0,0)$ is a technical point we ignore.
\end{heuristic}

This is a Schr\"odinger bridge, or entropic optimal transport, problem
\cite{Leonard}. The maximiser is $d\mu=a(x)a(y)\kappa(x,y)\,dx\,dy$, where
\[
a(x)\int_0^1 a(y)\kappa(x,y)\,dy=1\quad(0\le x\le1),\qquad\text{and}\qquad C=-2\int_0^1\log a(x)\,dx .
\]
Symmetric Sinkhorn iteration \cite{Sinkhorn} with the midpoint rule on $M=400,1600,3200$
nodes gives
$C=-0.632547,\ -0.632539,\ -0.632538$.

\begin{remark}[dictionary and references]\label{rem:eot}
Let $Z=\iint\kappa$ and $R=Z^{-1}\kappa\,dx\,dy$. Then
$C=\log Z-\min_{\mu\in\Pi(\lambda,\lambda)}\KL(\mu\,\|\,R)$. This is the static
Schr\"odinger bridge with cost $c=-\log\kappa$, regularisation $\varepsilon=1$ and uniform
marginals, in the notation of Nutz's notes \cite{Nutz}. The minimiser is
$e^{\varphi(x)+\psi(y)}R$ with Schr\"odinger potentials $e^\varphi=e^\psi=\sqrt Z\,a$, and
Sinkhorn's algorithm is alternating projection onto the two marginal constraints.
\begin{itemize}
\item Existence and uniqueness of the continuous scaling $a\,\kappa\,a$ (the ``DAD
problem'' for kernels): \cite{BLN,Nussbaum}.
\item Surveys: \cite{Idel} for matrix scaling, and \cite{Leonard,Nutz} for Schr\"odinger bridges.
\item Numerics: \cite{PeyreCuturi}.
\item The same variational formula as the limiting normalising constant of exponential
families on permutations, $\sum_\pi\exp(\theta\sum_i f(i/n,\pi(i)/n))$, computed by
iterative proportional fitting, with worked examples $f=|x-y|$ and $f=xy$:
\cite{Mukherjee,MukherjeeTagami}.
\item A fully worked example of the same entropy principle, for a two-body weight
(the Mallows model): \cite{Starr}.
\item Permanental means versus scaling means (cf.\ Lemma~L1 in \cref{sec:superexp}):
\cite{BochiIommiPonce}.
\end{itemize}
Most of these results assume a kernel bounded above and below. Here
$\kappa\ge1/(2\sqrt2)$, but $\kappa$ has an integrable singularity at $(0,0)$, so the
Hilbert-metric contraction argument of \cite{Nussbaum} does not apply verbatim, and
\cite{Mukherjee} needs a truncation of $\log\kappa$ near the corner.
\end{remark}

\begin{heuristic}[no arithmetic correction at exponential order]
$G(n)$ is not $R(n)$: its edges carry the local factors $\varrho_{\ell^e}(i+j)$. At
a prime power $q$, the local kernel $(r,s)\mapsto\varrho_q(r+s)$ on $\mathbb{Z}/q$ is
translation invariant, so it has uniform marginals. Hence the local analogue of
$C$, namely $\sup_\mu\{\int\log\varrho_q\,d\mu-\KL\}$, is attained at
$\mu\propto\varrho_q$ and equals $\log\int\varrho_q=0$. We therefore expect the
arithmetic (a ``singular series'') to affect $\log\Ham(n)$ only by $o(n)$.
\end{heuristic}

\begin{conjecture}\label{conj:asymp}
$\log\Ham(n)=\log n!-\frac n2\log n+Cn+o(n)$, with $C=-0.632538\ldots$ as above.
In particular $\Ham(n)=(\sqrt n/e)^{n}e^{Cn+o(n)}$.
\end{conjecture}

\begin{proposition}[rigorous upper bound]\label{prop:bregman}
Let $f(x)=\sqrt{1+x}-\sqrt x$ and
\[
B=\int_0^1\log f(x)\,dx=\frac{\sqrt2}{2}-\frac32\log(1+\sqrt2)=-0.6149536\ldots
\]
Then $\log\Ham(n)\le\frac n2\log n+(B-1)n+o(n)$. Moreover $C<B$.
\end{proposition}
\begin{proof}
Adjoin a vertex $0$ adjacent to all of $[n]$, and let $A'$ be the adjacency matrix of
the resulting graph. A directed Hamiltonian path $v_1\cdots v_n$ of $G(n)$ gives the
cycle cover $0\to v_1\to\cdots\to v_n\to0$ of the augmented graph. So
$2\Ham(n)\le\operatorname{per}A'\le\prod_i (r_i!)^{1/r_i}$ by Bregman's theorem
\cite{Bregman}, where $r_i=\deg_{G(n)}(i)+1$ for $i\ge1$ and $r_0=n$. Now
$\deg(i)=\sqrt n\,f(i/n)+O(1)$ uniformly, with $f\ge\sqrt2-1$, and
$\log (r!)^{1/r}=\log r-1+O(\log r/r)$. Summing gives the bound. For $B$, substitute
$x=\sinh^2t$: then $f=e^{-t}$, and $B=-\int_0^{\log(1+\sqrt2)}t\sinh 2t\,dt$.

For $C\le B$, disintegrate a coupling $\mu$ with first marginal $\lambda$ as
$\mu(dx\,dy)=dx\,\mu_x(dy)$. Then $\KL(\mu\|\lambda^2)=\int\KL(\mu_x\|\lambda)\,dx$.
By the Gibbs variational principle,
$\int\log\kappa(x,\cdot)\,d\mu_x-\KL(\mu_x\|\lambda)\le\log\int_0^1\kappa(x,y)\,dy=\log f(x)$.
Equality would force $\mu_x\propto\kappa(x,\cdot)$, which does not have a uniform second
marginal, so the inequality is strict.
\end{proof}
Since $\log n!-\frac n2\log n=\frac n2\log n-n+O(\log n)$, \cref{conj:asymp}
asserts that the true constant $C-1$ lies only $B-C\approx0.0175$ below the rigorous
bound $B-1$. The finite-$n$ Bregman bound converges to its limit slowly (at $n=3000$
the normalised value is still $\approx-0.49$), because the error terms are of order
$\log d/d$ with $d\asymp\sqrt n$.

\begin{table}[h]
\centering\small
\begin{tabular}{rrrrrrrrrrrr}
\toprule
$n$ & 15 & 25 & 30 & 35 & 40 & 45 & 50 & 55 & 60 & 65 & 70\\
$r(n)$ & $+1.24$ & $-0.57$ & $-2.68$ & $+0.89$ & $-1.81$ & $+0.10$ & $-1.18$ & $-0.87$ & $-1.81$ & $+0.34$ & $-1.82$\\
\midrule
$n$ & 71 & 72 & 73 & 74 & 75\\
$r(n)$ & $-1.18$ & $-1.50$ & $-1.12$ & $-0.75$ & $+3.92$\\
\bottomrule
\end{tabular}
\caption{Residuals $r(n)=\log\Ham(n)-\bigl[\log(n!/2)-\frac{n-1}2\log n+Cn\bigr]$
(natural logarithms) from \texttt{data/b071983.txt}. For $15\le n\le74$ (all $n$
with $\Ham(n)>0$), $r(n)\in[-2.68,1.24]$, while $\log\Ham(74)\approx42.2$. On
$40\le n\le74$ the mean is $-0.70$ and the standard deviation $0.72$. A least-squares
fit $r=a\log n+b$ gives $a\approx-1.0$, which is consistent with an extra $-\log n$
term, but the range is short.}
\label{tab:resid}
\end{table}

\begin{remark}
The value at $n=75$ is anomalous: $\Ham(75)/\Ham(74)\approx300$, whereas the ratios
$\Ham(n+1)/\Ham(n)$ for $70\le n\le73$ are between $2.0$ and $5.2$. The residual jumps by $4.7$.
Since vertex $75$ has degree $4$ in $G(75)$ (neighbours $6,25,46,69$), we suggest
rechecking $\Ham(75)$.
\end{remark}

\section{Towards superexponential lower bounds}\label{sec:superexp}

If \cref{conj:asymp} holds, the bounds of \cref{sec:moves} are far from the truth.
For random and pseudo-random graphs, the method of choice runs as follows.
(i) Bound the number of $2$-factors (cycle covers) from below via the permanent, using
the Egorychev--Falikman theorem \cite{Egorychev,Falikman} after Sinkhorn scaling; the
matching upper bound is Bregman's \cite{Bregman}. (ii) Show that covers with many
cycles contribute little. (iii) Merge the cycles into a Hamiltonian cycle by
P\'osa rotation--extension \cite{Posa}. This gives $n!p^n(1+o(1))^n$ Hamiltonian cycles in
$G(n,p)$ above the Hamiltonicity threshold \cite{GlebovKrivelevich} and in
$(n,d,\lambda)$-graphs \cite{Krivelevich}. See \cite{CooperFrieze,CucklerKahn,FKS} for
related counts.

For $G(n)$, step (i) should give, rigorously, about $(\sqrt n/e)^n e^{Cn}$ cycle
covers, provided the Sinkhorn scalings of $A_{G(n)}$ track the continuum solution $a(x)$.
The arithmetic enters only through $o(n)$ terms, by the heuristic above. Step (iii)
needs expansion of small vertex sets in $G(n)$. Here the arithmetic is a genuine obstacle: $G(n)$ is far from
pseudo-random in the usual sense. Squares are $0$ or $1$ modulo $4$, so the residue
classes $1$ and $3$ modulo $4$ are independent sets of size about $n/4$, and there are no
edges between the classes $0$ and $2$ modulo $4$. Likewise, the class $1$ modulo $3$ is
an independent set of size about $n/3$. Every allowed pair of classes mod $4$ has twice the
$\kappa$-density. Rotation--extension arguments that rely on edges between any two
large sets therefore have to be made residue-aware. This is where rotations would do
real work.
By contrast, a single rotation changes an endpoint. In the lifting scheme the
endpoints of each $T_c$ are pinned by the glue, so rotations contribute only constant
factors per level unless they are composed into the fixed-endpoint moves of
\cref{def:move}.

A different general mechanism, rerouting inside consecutive blocks of a
Hamiltonian path, gives exponential counts for traceable graphs with certain forbidden
induced subgraphs \cite{JookenZamfirescu}. It requires every sufficiently long block
to be non-rigid. In $G(n)$ chords between path vertices at bounded distance become rare
as $n\to\infty$, so it does not apply directly. That is why \cref{sec:moves} needs
the lift.

\section{A deterministic route for $G(n)$}\label{sec:deterministic}

The random graph $R(n)$ of \cref{heu:first} was used only to \emph{guess} the constant
$C$. \Cref{conj:asymp} is a statement about the deterministic family $G(n)$. In the
proofs for random graphs \cite{GlebovKrivelevich}, randomness is used only to verify a
few structural inputs. Krivelevich's version for pseudo-random graphs \cite{Krivelevich}
makes these inputs explicit and is entirely deterministic. In this section we isolate the
corresponding statements for $G(n)$. We show that the first of them, the count of cycle
covers, already transfers to $G(n)$ in a deterministic way, and we verify it numerically.
\Cref{app:limits} explains the analytic background in more detail.

Throughout, $A=A_{G(n)}$ is the adjacency matrix. Then $\operatorname{per}A$ is the number
of permutations $\pi$ of $[n]$ with $i\sim\pi(i)$ for all $i$. These are the cycle covers
of the directed graph obtained by doubling each edge; $2$-cycles are allowed.
Directed Hamiltonian cycles are among them, so $2\cdot\#\{\text{Hamiltonian cycles}\}\le\operatorname{per}A$.

\subsection{Cycle covers: a deterministic lower bound}

\begin{theorem}[Gurvits \cite{Gurvits}; Gurvits--Samorodnitsky {\cite[Thm.~1.2]{GurvitsSamorodnitsky}}]\label{thm:GS}
Let $A$ be a nonnegative $n\times n$ matrix and $P$ a doubly stochastic matrix. Then
\[
\log\operatorname{per}A\ \ge\ \sum_{i,j}\Bigl[p_{ij}\log\frac{a_{ij}}{p_{ij}}+(1-p_{ij})\log(1-p_{ij})\Bigr],
\]
with the conventions $0\log(\cdot/0)=0$ and $\log 0=-\infty$.
\end{theorem}

The maximum of the right-hand side over $P$ is the logarithm of the \emph{Bethe
permanent} of $A$. The theorem is equivalent to Gurvits's inequality
$\operatorname{per}A\ge\operatorname{Bethe}(A)$, which he derived from Schrijver's
permanental inequality \cite{Schrijver}. For a $0/1$ matrix and $P$ supported on the ones
of $A$, the right-hand side is
\[
F(P)=\sum_{i\sim j}\bigl[-p_{ij}\log p_{ij}+(1-p_{ij})\log(1-p_{ij})\bigr].
\]

\begin{corollary}\label{cor:bethe}
For every doubly stochastic $P$ supported on the edges of $G(n)$ (each ordered pair
counted once), $\log\operatorname{per}A_{G(n)}\ge F(P)$.
\end{corollary}

Nothing here is random. The only question is how good a $P$ we can exhibit. Take $P=S$,
the symmetric Sinkhorn scaling $S=DAD$ of $A$, computed numerically to machine precision.
This gives:

\begin{center}
\begin{tabular}{rcc}
\toprule
$n$ & $\bigl(F(S)-\frac n2\log n\bigr)/n$ & max.\ row-sum error of the ansatz \eqref{eq:ansatz}\\
\midrule
$1000$ & $-1.60375$ & $0.107$\ (rms $0.026$)\\
$4000$ & $-1.61752$ & $0.067$\ (rms $0.013$)\\
$16000$ & $-1.62498$ & $0.031$\ (rms $0.006$)\\
$64000$ & $-1.62871$ & ---\\
\midrule
$\infty$ (continuum) & $C-1=-1.63254$ &\\
\bottomrule
\end{tabular}
\end{center}

The gap to $C-1$ halves each time $n$ is multiplied by $4$, so it decays like $n^{-1/2}$.
Numerically, therefore,
$\log\operatorname{per}A_{G(n)}\ge\frac n2\log n+(C-1)n-O(\sqrt n)$. This is a
deterministic check that the residue structure of $G(n)$ does not change the constant
at the level of cycle covers. In the other direction, Bregman's theorem \cite{Bregman}
gives the proved bound
\[
\log\operatorname{per}A_{G(n)}\le\sum_i\frac{\log(d_i!)}{d_i}=\frac n2\log n+(B-1)n+o(n),
\]
exactly as in \cref{prop:bregman}.

\subsection{Why the discrete scaling tracks the continuum}
The scaling of $A$ is close to the continuum solution $a$ of \cref{sec:heuristic}
because its row sums are sums of a fixed function over the squares in an interval. No
congruence conditions enter.

\begin{lemma}\label{lem:squaresum}
Let $f$ have bounded variation on $[0,1]$, let $\delta>0$, and let $1\le i\le n$ with
$x=i/n\ge\delta$. Then
\[
\sum_{k:\ i<k^2\le i+n} f\Bigl(\frac{k^2-i}{n}\Bigr)
=\sqrt n\int_0^1 f(y)\,\kappa(x,y)\,dy+O_\delta\bigl(\|f\|_\infty+V(f)\bigr).
\]
\end{lemma}

\begin{proof}
Put $y_k=(k^2-i)/n$. Then $y_{k+1}-y_k=(2k+1)/n\le 3/\sqrt n$, and since
$x+y_k=k^2/n$,
\[
\sqrt n\,\kappa(x,y_k)=\frac{n}{2k},\qquad
f(y_k)=f(y_k)\,\sqrt n\,\kappa(x,y_k)\,(y_{k+1}-y_k)\Bigl(1+O\bigl(\tfrac1k\bigr)\Bigr).
\]
So the sum is a Riemann sum for $\int_0^1 f(y)\,\sqrt n\,\kappa(x,y)\,dy$. The weight $\sqrt n\,\kappa(x,\cdot)$ is
monotone and bounded by $\sqrt n/(2\sqrt\delta)$. The mesh is $O(n^{-1/2})$, and there are
$O(\sqrt n)$ terms each with relative error $O(1/k)=O(n^{-1/2})$. So the total error is
$O_\delta(\|f\|_\infty+V(f))$.
\end{proof}

Applying this with $f=a$ (or rather a bounded-variation truncation of it near $0$) shows
that the \emph{ansatz}
\begin{equation}\label{eq:ansatz}
\widetilde p_{ij}=\frac{a(i/n)\,a(j/n)}{\sqrt n}\qquad(i\sim j)
\end{equation}
has row and column sums $1+O_\delta(n^{-1/2})$ on rows with $i\ge\delta n$. This uses
$a(x)\int_0^1 a\,\kappa(x,\cdot)=1$. The table above measures this error directly.
Moreover (\cref{app:bethe}),
$F(\widetilde p)=\frac n2\log n+(C-1)n+o(n)$. What remains is to correct $\widetilde p$
to an exactly doubly stochastic matrix:

\begin{problem}[L1]\label{prob:L1}
Show that there is a doubly stochastic $P$ supported on $E(G(n))$ with
$\sum_{ij}|p_{ij}-\widetilde p_{ij}|\,\log n=o(n)$. With \cref{cor:bethe} this would give
\[
\log\operatorname{per}A_{G(n)}\ \ge\ \frac n2\log n+(C-1)n-o(n),
\]
the lower bound of \cref{conj:asymp} for cycle covers (the upper bound being Bregman's,
with $B$ in place of $C$).
\end{problem}

This is a flow problem on the bipartite double cover of $G(n)$: route the excess
$O(n^{-1/2})$ per row along short paths. The quotient graphs on residue classes are
connected (see \cref{subsec:merge}), so there is no arithmetic obstruction. Rows with
$i<\delta n$ need separate treatment because of the singularity of $\kappa$ at the origin.

\subsection{Few cycles}
\begin{problem}[L2]\label{prob:L2}
Show that for some $s=s(n)=o(n/\log n)$ the cycle covers of $G(n)$ with at most $s$
cycles number at least $e^{-o(n)}\operatorname{per}A_{G(n)}$.
\end{problem}
In \cite{Krivelevich} this is done deterministically. A cover with at least $s$ cycles has
at least $s/2$ cycles of length at most $2n/s$. One bounds the covers containing a
prescribed family of short cycles by applying Bregman to the remaining matrix, and then
sums over families. So the input is an upper bound on the number $N_L(n)$ of
$L$-cycles of $G(n)$. Such a cycle $x_1\cdots x_L$ is a solution of
$x_t+x_{t+1}=a_t^2$ ($t$ mod $L$). Eliminating the $x_t$, this means
$\sum_t(-1)^ta_t^2=0$ for even $L$, and $2x_1=\sum_t(-1)^{t+1}a_t^2$ for odd $L$, together with
ordering constraints. Counting these is classical Diophantine analysis, in which
arithmetic enters only through logarithmic factors. (For $L=4$, $a^2+c^2=b^2+d^2$ has
$\asymp N^2\log N$ solutions with $a,b,c,d\le N$.) A spectral variant: under the
Gibbs measure on covers, the expected number of cycles is heuristically
$\sum_L\operatorname{tr}(S^L)/L$. Only finitely many eigenvalues of $S$, coming from
characters of small moduli, are of constant size, which suggests $O(\sqrt n)$ cycles.

\subsection{Merging}\label{subsec:merge}
\begin{problem}[L3]\label{prob:L3}
Prove a residue-aware expansion property of $G(n)$ strong enough for the
rotation--extension merging of \cite{Krivelevich,GlebovKrivelevich}.
\end{problem}
Here the arithmetic is essential. The class $1\bmod 3$ is independent, and there are no
edges between the classes $0$ and $2\bmod 4$. The quotient graphs are nevertheless
connected: modulo $4$ the class pairs that occur are
$\{0,0\},\{0,1\},\{1,3\},\{2,2\},\{2,3\}$, and modulo $3$ they are
$\{0,0\},\{0,1\},\{1,2\},\{2,2\}$. So a natural formulation is expansion
\emph{relative to} the residue structure.

\subsection{The upper bound}
The deterministic upper bound (Bregman) has constant $B$, not $C$. Universal comparisons
between the permanent and its Bethe approximation lose a factor exponential in $n$:
$\operatorname{per}A\le 2^{n/2}\operatorname{Bethe}(A)$ \cite{AnariRezaei}, with
improvements to $c^n$, $c<\sqrt2$, cited in \cite{DongJain}. Girth-dependent
improvements \cite{DongJain} do not apply, since $G(n)$ contains triangles
($6+19=25$, $19+30=49$, $30+6=36$). A sharp upper bound would need a ``large-degree
Bethe-exactness'' theorem for deterministic matrices with all scaled entries $o(1)$.
We do not know such a result. The realistic target is therefore the lower bound with
constant $C$ together with the upper bound with constant $B$, a gap of $0.0175$ per vertex.

\subsection{Barvinok's identity and a finite-$n$ predictor}\label{subsec:barvinok}
Barvinok \cite[Lemma~4.2]{Barvinok} defines the \emph{Hamiltonian permanent}
$\operatorname{hamp}A=\sum_{g}\prod_ia_{i,g(i)}$, the sum over permutations $g$ consisting of a
single $n$-cycle, and proves
$\operatorname{hamp}A=\frac1n[x_1\cdots x_n]\operatorname{Tr}(D(x)A)^n$ with
$D(x)=\operatorname{diag}(x)$. The same argument for paths gives
\[
2\Ham(n)=[x_1\cdots x_n]\ \mathbf 1^{\!\top}(D(x)A)^{n-1}D(x)\mathbf 1,\qquad A=A_{G(n)} .
\]
Unlike $\operatorname{per}A=[x_1\cdots x_n]\prod_i(Ax)_i$, this coefficient sees only single
paths. In principle it therefore bypasses \cref{prob:L2,prob:L3}.

\emph{An exact identity.} Let $S=DAD$ be the symmetric Sinkhorn scaling, $D=\operatorname{diag}(d_j)$,
and evaluate at $x_j=d_j^2$. A walk $v_1\cdots v_n$ then has weight
$\prod_t d_{v_t}^2\prod_tA_{v_tv_{t+1}}=d_{v_1}d_{v_n}\prod_tS_{v_tv_{t+1}}$. Let $(X_t)$ be the stationary
Markov chain with transition matrix $S$; its stationary law is uniform because $S$ is doubly
stochastic. Then
\begin{equation}\label{eq:markov}
2\Ham(n)=n\Bigl(\prod_jd_j\Bigr)^{-2}\,\mathbb{E}\bigl[d_{X_1}d_{X_n};\ X_1,\dots,X_n\text{ distinct}\bigr].
\end{equation}

\emph{Zeroth order.} For i.i.d.\ uniform $X_t$, $\mathbb{P}(\text{distinct})=n!/n^n$.
Using this value and decoupling the endpoints gives the predictor
\[
\widehat\Ham(n)=\frac n2\cdot\frac{n!}{n^n}\cdot\Bigl(\prod_jd_j\Bigr)^{-2}\Bigl(\frac1n\sum_jd_j\Bigr)^{2}.
\]
If $d_j\approx n^{-1/4}a(j/n)$ (\cref{prob:L1}), then
$\log\widehat\Ham(n)=\log n!-\frac n2\log n+Cn+o(n)$, which is \cref{conj:asymp}. Unlike the formula
of \cref{tab:resid}, $\widehat\Ham(n)$ uses the actual scaling of each $G(n)$.

\emph{What fails.} Two natural refinements do not work at the sizes where we can test them.
\begin{enumerate}
\item A Gaussian saddle-point approximation of the torus integral is wrong at exponential order.
The visit counts are lattice variables with mean $1$ and variance $\approx1$. Their local
probability is $\approx e^{-1}$ (Poisson), not the Gaussian $(2\pi)^{-1/2}$. Already for
$S=J/n$ this gives $(2\pi)^{-(n-1)/2}$ in place of $n!/n^n$, an error of about $0.081n$ in the
logarithm. The correct baseline is the i.i.d.\ one used above.
\item First-order ``loop'' corrections to the i.i.d.\ baseline are far too large.
These are the pair-collision factor $\exp\bigl(-\sum_\ell(1-\ell/n)(\operatorname{tr}S^\ell-1)\bigr)$
and the cycle-cover analogue $\det'(I-S)=\prod_{\lambda\ne1}(1-\lambda)$. For $31\le n\le75$ they
change $\log\widehat\Ham$ by $-23$ to $-29$ and by $-8$ to $-10$ respectively, whereas the
required correction is about $+2$. At these sizes the degrees are only $2$--$7$,
$\|S\|_F^2-1\approx12$--$17$, and $S$ has eigenvalues down to $-0.95$, so the chain is far
from the perturbative regime.
\end{enumerate}

\emph{What holds.} The zeroth-order predictor, by contrast, works remarkably well; see
\cref{fig:predictor}. For $31\le n\le74$,
$\log\Ham(n)-\log\widehat\Ham(n)$ lies in $[1.89,3.11]$, with standard deviation $0.23$. The
constant-$C$ formula has standard deviation $0.72$. The new residual varies smoothly, changing
by at most $0.6$ between consecutive $n$, and usually by less than $0.15$. In particular, the
oscillations of the residual in \cref{tab:resid}, with dips near $n=40,58$--$60,70$, are
explained entirely by the scaling vector, that is, by the capacity of $A_{G(n)}$. What remains
is a slowly varying correction of size $\approx2$, consistent with $e^{o(n)}$. A least-squares
fit $r\approx1.0+8.9/\sqrt n$ is empirical only. For $n\le30$ the Sinkhorn iteration
degenerates (some entries of $S$ tend to $0$), because $A_{G(n)}$ lacks total support there.

\begin{figure}[h]
\centering
\includegraphics[width=0.85\linewidth]{../figs/predictor.pdf}
\caption{Residuals of $\log\Ham(n)$ against the constant-$C$ formula of \cref{tab:resid} and
against the scaling predictor $\widehat\Ham(n)$, for $31\le n\le75$. The point at $n=75$ is
the suspect value.}
\label{fig:predictor}
\end{figure}

\emph{The value at $n=75$.} Against $\widehat\Ham$ the residual at $n=75$ is $+6.25$, a jump of
$+4.3$ from a curve that otherwise moves by at most $0.6$ per step. This strengthens the case
for rechecking $\Ham(75)$.

In terms of \eqref{eq:markov}, everything beyond \cref{prob:L1} reduces to one probabilistic
statement:
\begin{problem}[L$^*$]\label{prob:Lstar}
For the stationary chain with transition matrix $S$, the Sinkhorn scaling of $A_{G(n)}$, show that
$\log\mathbb{P}(X_1,\dots,X_n\text{ distinct})=\log(n!/n^n)+o(n)$.
\end{problem}
Together with \cref{prob:L1}, a solution would prove \cref{conj:asymp} in full, for paths
directly, without passing through cycle covers. The data suggest that the ratio to $n!/n^n$ is
in fact $e^{O(1)}$ in this range.


\section{Software}\sloppy
The code is in \texttt{hamiltonian/bounds}; run it with
\texttt{uv run python -m hamiltonian.bounds.<module>}.
\texttt{certificate} verifies \cref{cor:seeds} for $m\le2$ from
\texttt{data/bounds/certificate.json}. \texttt{glue} gives \cref{prop:t3}
(\texttt{--t 3 --exhaustive}). \texttt{seeds} searches for seeds. \texttt{sinkhorn}
computes $C$. \texttt{residuals} produces \cref{tab:resid}. \texttt{edges} checks
\cref{prop:edges,prop:mod4}. \texttt{bregman} computes the bound of \cref{prop:bregman},
\texttt{bethe} produces the table in \cref{sec:deterministic}, and \texttt{hamp\_predictor}
produces \cref{fig:predictor}.

\subsection*{Acknowledgement}
The arguments and computations in this draft were developed in conversation with
Claude (Anthropic), and should be checked independently.

\appendix
\section{Graph limits, permutons, and the Schr\"odinger problem}\label{app:limits}

This appendix explains in more detail the analytic objects used in
\cref{sec:limits,sec:heuristic,sec:deterministic}. Its aim is to let a reader who has not
met graphons or entropic transport check each step. Statements marked as theorems are
quoted from the literature or proved here. Heuristic steps are labelled as such.

\subsection{Dense graph limits}
For a graph $G$ on $[n]$, define the step function $W_G$ on $[0,1]^2$ to be $1$ on
$\bigl(\frac{i-1}n,\frac in\bigr]\times\bigl(\frac{j-1}n,\frac jn\bigr]$ when $i\sim j$ and
$0$ otherwise. Thus $W_G$ is ``the adjacency matrix drawn as a picture''. A
\emph{graphon} is a symmetric measurable $W:[0,1]^2\to[0,1]$. For a fixed graph $F$
on $[k]$ the homomorphism density is
\[
t(F,W)=\int_{[0,1]^k}\prod_{\{u,v\}\in E(F)}W(x_u,x_v)\,dx_1\cdots dx_k .
\]
$t(F,W_G)$ is the probability that a random map $V(F)\to V(G)$ is a homomorphism. The
\emph{cut norm} is $\|W\|_\square=\sup_{S,T}\bigl|\int_{S\times T}W\bigr|$, the supremum
over measurable $S,T\subseteq[0,1]$. The \emph{cut distance}
$\delta_\square(W,U)=\inf_\phi\|W-U^\phi\|_\square$ takes the infimum over
measure-preserving relabellings $\phi$ of $[0,1]$.

\begin{theorem}[Lov\'asz--Szegedy \cite{LovaszSzegedy}; see \cite{Lovasz}]
Every sequence of graphs with $|V(G_n)|\to\infty$ has a subsequence along which
$t(F,W_{G_n})$ converges for every $F$. The limits are $t(F,W)$ for a graphon $W$, and
convergence of all $t(F,\cdot)$ is equivalent to $\delta_\square(W_{G_n},W)\to0$.
\end{theorem}

Conversely, a graphon $W$ generates \emph{$W$-random graphs}: choose
$x_1,\dots,x_n$ independently and uniformly from $[0,1]$, and join $i,j$ independently with
probability $W(x_i,x_j)$. These converge to $W$ almost surely. For example,
$G(n,\tfrac12)\to\tfrac12$, and the graphs with $i\sim j\iff i+j\le n$ converge to
$\mathbf 1[x+y\le1]$.

\subsection{Sparse graphs, and why $G(n)$ needs more}
If $e(G_n)=o(n^2)$ then $W_{G_n}\to0$, which carries no information. The theory of $L^p$
graphons \cite{BCCZ1,BCCZ2} rescales by the edge density $\rho_n=2e(G_n)/n^2$. One
studies $W_{G_n}/\rho_n$, whose limits are symmetric $W\in L^p([0,1]^2)$, $p>1$, with no
upper bound required. Compactness requires \emph{$L^p$ upper regularity}: block averages of
$W_{G_n}/\rho_n$ over partitions into parts of measure at least $\eta$ must have $L^p$ norm
at most a constant $K$ independent of $\eta$. The sampling model joins $i,j$ with
probability $\min(1,\rho_nW(x_i,x_j))$.

For $G(n)$ we have $\rho_n\asymp n^{-1/2}$. The natural candidate limit is $\kappa(x,y)=1/(2\sqrt{x+y})$:
the squares near $m$ have density $1/(2\sqrt m)$, and
$\sqrt n\cdot 1/(2\sqrt{n(x+y)})=\kappa(x,y)$. \Cref{prop:mod4} shows that cut-norm
convergence to $\kappa$ fails, and the remark after it shows that upper regularity fails.
Both failures come from residues: whether $i+j$ is a square depends on $i\bmod q$ and
$j\bmod q$ for every $q$. The limit must therefore record residues.

\subsection{The profinite integers}
$\Zhat=\varprojlim\mathbb{Z}/q\mathbb{Z}\cong\prod_p\mathbb{Z}_p$ is a compact ring
containing $\mathbb{Z}$ densely. An element is a compatible choice of residues modulo every
$q$. Its Haar probability measure $\mu_{\Zhat}$ gives each residue class modulo $q$ mass
$1/q$. Let $\sigma$ be the law of $k^2$ when $k$ is Haar-distributed. With
$\varrho_q(a)=\#\{x\bmod q:\ x^2\equiv a\}$,
\[
\sigma(\{s\equiv a\bmod q\})=\varrho_q(a)/q .
\]
The squares in $\mathbb{Z}_p$ have Haar measure $\sum_{k\ge0}p^{-2k}\cdot\frac{p-1}{2p}=\frac{p}{2(p+1)}$
for odd $p$, and $\sum_{k\ge0}4^{-k}\cdot\frac18=\frac16$ for $p=2$. The infinite
product over $p$ is $0$, so the squares form a Haar-null set in $\Zhat$, and $\sigma$
is singular with respect to $\mu_{\Zhat}$.

\begin{proposition}[joint equidistribution]\label{prop:equid}
Fix $q$, residues $r,s\bmod q$, and intervals $I,J\subseteq[\delta,1]$, $\delta>0$. Then
\[
\#\bigl\{(i,j):\ i\sim j,\ \tfrac in\in I,\ \tfrac jn\in J,\ i\equiv r,\ j\equiv s\ (\mathrm{mod}\ q)\bigr\}
=\frac{\varrho_q(r+s)}{q^2}\,n^{3/2}\!\iint_{I\times J}\kappa+O_{q,\delta}(n).
\]
\end{proposition}

\begin{proof}
Fix $i\equiv r$. Its neighbours are $j=k^2-i$, and $j\equiv s$ if and only if
$k^2\equiv r+s\pmod q$. Within any $q$ consecutive values of $k$ there are exactly
$\varrho_q(r+s)$ such $k$. Apply \cref{lem:squaresum} with $f=\mathbf 1_J$ to the
subsequence of $k$ in a fixed residue class; this changes the mesh by the factor $q$.
The count for fixed $i$ is $\frac{\varrho_q(r+s)}q\sqrt n\int_J\kappa(i/n,y)\,dy+O_{q,\delta}(1)$.
Now sum over the $\frac nq|I|+O(1)$ values of $i$.
\end{proof}

Accordingly, put $\Omega=[0,1]\times\Zhat$ with measure $\lambda\otimes\mu_{\Zhat}$, and send
vertex $i$ to $(i/n,i)$. Define the operator
\[
(Tf)(x,\alpha)=\int_0^1\kappa(x,y)\,\mathbb{E}_{s\sim\sigma}\,f(y,\,s-\alpha)\,dy .
\]
\Cref{prop:equid} says exactly that the normalised adjacency operator
$(A_nf)(i)=n^{-1/2}\sum_{j\sim i}f(j)$ converges to $T$ against all test functions that
depend on $\alpha$ only through a fixed modulus. In the language of Backhausz--Szegedy
\cite{BackhauszSzegedy}, a \emph{graphop} is a positivity-preserving self-adjoint
operator $L^\infty(\Omega)\to L^1(\Omega)$. \emph{Action convergence} asks that, for
every $k$ and all $f_1,\dots,f_k$ with $\|f_i\|_\infty\le1$, the joint laws of
the $2k$-tuples
\[(f_1,\dots,f_k,\,A_nf_1,\dots,A_nf_k)\]
be approximated by those for $T$, and conversely.
\Cref{conj:limit} asserts this for $T$. None of the arguments below depends on it. Its
role is to explain why a single analytic object, $\kappa$ together with $\sigma$,
should govern all of the asymptotics.

\subsection{Permutons and their large deviations}
A \emph{permuton} is a probability measure on $[0,1]^2$ with both marginals uniform; the
set of permutons is $\Pi(\lambda,\lambda)$. A permutation $\pi\in S_n$ defines the
permuton $\mu_\pi$ with density $n$ on the squares
$\bigl(\frac{i-1}n,\frac in\bigr]\times\bigl(\frac{\pi(i)-1}n,\frac{\pi(i)}n\bigr]$. We use the weak
topology on permutons.

\begin{theorem}[Trashorras; Kenyon--Kr\'al'--Radin--Winkler {\cite[Thm.~1]{KKRW}}]\label{thm:LDPperm}
If $\pi$ is uniform in $S_n$, the laws of $\mu_\pi$ satisfy a large deviation principle
with speed $n$ and rate function $I(\mu)=\KL(\mu\,\|\,\lambda^2)=\int\rho\log\rho$ if $\mu$
has density $\rho$, and $I=+\infty$ otherwise. Informally,
$\#\{\pi:\mu_\pi\approx\mu\}\approx n!\,e^{-n\KL(\mu\|\lambda^2)}$.
\end{theorem}

By Varadhan's lemma \cite[Thm.~4.3.1]{DemboZeitouni}, for bounded continuous $\Phi$ on
permutons,
\[
\lim_{n\to\infty}\frac1n\log\mathbb{E}_\pi\,e^{n\Phi(\mu_\pi)}=\sup_\mu\{\Phi(\mu)-I(\mu)\}.
\]
We apply this to $\Phi(\mu)=\int\log\kappa\,d\mu$. Since $\kappa\ge 1/(2\sqrt2)$,
$\log\kappa$ is bounded below. It is unbounded above only at the corner, where
$\kappa(x,y)\le 1/(2\sqrt x)$. Let $\kappa_M=\min(\kappa,M)$. Only rows with
$i<n/(4M^2)$ can have $\kappa>M$, so for every permutation
\[
0\le\sum_i\Bigl[\log\kappa\bigl(\tfrac in,\tfrac{\pi(i)}n\bigr)-\log\kappa_M\bigl(\tfrac in,\tfrac{\pi(i)}n\bigr)\Bigr]
\le\sum_{i<n/(4M^2)}\tfrac12\log\frac{n}{4M^2i}=O\!\left(\frac n{M^2}\right).
\]
Also $\sum_i\log\kappa_M(i/n,\pi(i)/n)=n\int\log\kappa_M\,d\mu_\pi+o(n)$ uniformly in $\pi$,
by uniform continuity. Letting $M\to\infty$ after $n\to\infty$ proves:

\begin{proposition}\label{prop:firstmoment}
Let $R_n$ be the random $0/1$ matrix with independent entries
$\mathbb{P}(R_n(i,j)=1)=p_{ij}=\frac1{2\sqrt{i+j}}$. Then
\[
\mathbb{E}\operatorname{per}R_n=\sum_{\pi\in S_n}\prod_ip_{i\pi(i)}=\frac{n!}{n^{n/2}}\ \mathbb{E}_\pi\prod_i\kappa\bigl(\tfrac in,\tfrac{\pi(i)}n\bigr),
\qquad
\frac1n\log\mathbb{E}_\pi\prod_i\kappa\bigl(\tfrac in,\tfrac{\pi(i)}n\bigr)\longrightarrow C,
\]
where $C=\sup_{\mu\in\Pi(\lambda,\lambda)}\{\int\log\kappa\,d\mu-\KL(\mu\|\lambda^2)\}$.
\end{proposition}

This is a theorem, given \cref{thm:LDPperm}. The \emph{heuristic} steps of
\cref{sec:heuristic} are: (i) the passage from $\mathbb{E}\operatorname{per}R_n$, the
first moment of the number of cycle covers of the random model, to Hamiltonian cycles,
i.e.\ to cyclic $\pi$, at a cost of $e^{o(n)}$; and (ii) the passage from $R_n$ to $G(n)$.
\Cref{sec:deterministic} replaces (ii) for cycle covers by a deterministic argument.

\subsection{The Schr\"odinger problem}
We now analyse the variational problem
\[
C=\sup_{\mu\in\Pi(\lambda,\lambda)}\Bigl\{\int\log\kappa\,d\mu-\KL(\mu\,\|\,\lambda^2)\Bigr\}.
\]

\emph{Gibbs variational principle.} For a probability measure $\nu$ and measurable $g$
with $\int e^g\,d\nu<\infty$,
\[
\sup_\mu\Bigl\{\int g\,d\mu-\KL(\mu\|\nu)\Bigr\}=\log\int e^g\,d\nu,
\]
attained exactly at $d\mu\propto e^g\,d\nu$. This follows from Jensen's inequality applied to
$\log\int e^{g}\frac{d\nu}{d\mu}\,d\mu$. Without the marginal constraints, our
supremum would therefore be $\log\iint\kappa=\log\frac43(\sqrt2-1)$. The constraints are what
make the problem a \emph{Schr\"odinger bridge}: with
$R=\kappa\,dx\,dy/\iint\kappa$,
\[
C=\log\iint\kappa-\min_{\mu\in\Pi(\lambda,\lambda)}\KL(\mu\,\|\,R).
\]
This is the static Schr\"odinger problem, equivalently entropic optimal transport with cost
$c=-\log\kappa$ and regularisation parameter $1$ \cite{Leonard,Nutz,PeyreCuturi}.

\emph{Existence and uniqueness.} $\Pi(\lambda,\lambda)$ is weakly compact, $\KL(\cdot\|R)$
is lower semicontinuous and strictly convex, and $\KL(\mu\|R)<\infty$ for $\mu=\lambda^2$.
So there is a unique minimiser $\mu^*$.

\emph{Form of the minimiser.} Perturb $\mu$ within $\Pi(\lambda,\lambda)$ and introduce
Lagrange multipliers $\varphi(x),\psi(y)$ for the two marginal constraints. Stationarity of
$\int\log\kappa\,d\mu-\int\rho\log\rho$ gives $\log\kappa-\log\rho-1=\varphi(x)+\psi(y)$,
that is,
\[
d\mu^*=a(x)\,b(y)\,\kappa(x,y)\,dx\,dy .
\]
The rigorous version, the product form of $I$-projections onto sets of measures with
given marginals, is due to Csisz\'ar \cite{Csiszar}. Borwein--Lewis--Nussbaum
\cite{BLN} treat kernels in exactly this ``$DAD$'' form, including measurable kernels. Since
$\kappa$ is symmetric and $\mu^*$ is unique, $\mu^*$ is symmetric. Then $a(x)b(y)=a(y)b(x)$,
so $b$ is a constant multiple of $a$, and after rescaling $b=a$. The marginal constraints
become the \emph{Schr\"odinger equation}
\begin{equation}\label{eq:schrod}
a(x)\int_0^1a(y)\,\kappa(x,y)\,dy=1\qquad(0\le x\le1).
\end{equation}

\emph{The value.} Since $\log\frac{d\mu^*}{d\lambda^2}=\log a(x)+\log a(y)+\log\kappa$, we get
$\int\log\kappa\,d\mu^*-\KL(\mu^*\|\lambda^2)=-\int(\log a(x)+\log a(y))\,d\mu^*=-2\int_0^1\log a$,
using the uniform marginals. Hence $C=-2\int_0^1\log a$.

\emph{A dual formula giving certified bounds.} For any measurable $u>0$ and any
$\mu\in\Pi(\lambda,\lambda)$, apply the Gibbs principle with $g=\log(\kappa\,u\otimes u)$, and use
$\int\log(u\otimes u)\,d\mu=2\int\log u$:
\[
\int\log\kappa\,d\mu-\KL(\mu\|\lambda^2)\ \le\ \log\iint\kappa(x,y)u(x)u(y)\,dx\,dy-2\int_0^1\log u .
\]
Equality holds for $u=a$ by \eqref{eq:schrod}. So
\[
C=\min_{u>0}\Bigl\{\log\iint\kappa\,u\otimes u-2\int\log u\Bigr\}
 =\max_{\mu\in\Pi(\lambda,\lambda)}\Bigl\{\int\log\kappa\,d\mu-\KL(\mu\|\lambda^2)\Bigr\}.
\]
Any trial function $u$ gives an upper bound for $C$, and any permuton gives a lower bound.
With $u=1$ we get $C\le\log\iint\kappa=-0.5937$. The comparison $C<B$ of
\cref{prop:bregman} is the same principle with only the first marginal constrained.

\subsection{Sinkhorn's iteration}
Write $(Kf)(x)=\int_0^1\kappa(x,y)f(y)\,dy$. Sinkhorn's iteration
\cite{Sinkhorn,SinkhornKnopp} alternately enforces the two marginals:
$b_{k+1}=1/(Ka_k)$, $a_{k+1}=1/(Kb_{k+1})$. A fixed point satisfies $a\,Kb=1=b\,Ka$. In the
symmetric case we use the damped form
\[
a_{k+1}=\bigl(a_k/Ka_k\bigr)^{1/2},
\]
the geometric mean of $a_k$ and the Sinkhorn update. Its fixed points are exactly the
solutions of \eqref{eq:schrod}.

\emph{Convergence.} On positive functions, use Hilbert's projective metric
$d_H(f,g)=\log\bigl(\sup\frac fg\cdot\sup\frac gf\bigr)$. The map $f\mapsto1/f$ is a
$d_H$-isometry. By the Birkhoff--Hopf theorem \cite{Birkhoff}, a positive linear operator
whose image has projective diameter $\Delta<\infty$ contracts $d_H$ by the factor
$\tanh(\Delta/4)$. Taking geometric means halves distances. So the damped iteration
contracts with factor at most $\frac12\bigl(1+\tanh(\Delta/4)\bigr)<1$ whenever
$\kappa$ is bounded above and below, and then $\Delta\le2\log(\sup\kappa/\inf\kappa)$.
Our $\kappa$ is unbounded only at the origin. For the truncation $\kappa_M$ the iteration
converges geometrically. $C$ is monotone in the kernel, and the estimate preceding
\cref{prop:firstmoment} gives $0\le C-C_M=O(M^{-2})$. Nussbaum \cite{Nussbaum} treats
the scaling of integral kernels in this Hilbert-metric framework in much greater generality.

\emph{Numerics.} With the midpoint rule on $M$ nodes, $K$ becomes an $M\times M$ matrix.
For $M=400,1600,3200$ the iteration converges to residual below $10^{-12}$ and gives
$C=-0.632547,\,-0.632539,\,-0.632538$ (\texttt{hamiltonian.bounds.sinkhorn}). The potential
$a$ increases from about $0.82$ at $x=0$, a finite value because $\kappa(0,\cdot)$ is
integrable, to about $1.76$ at $x=1$. It is plotted in the talk figures.

\subsection{The discrete counterpart: scaling, permanents, Bethe}\label{app:bethe}
The finite-$n$ analogue of a permuton is a doubly stochastic matrix. The analogue of the
Schr\"odinger problem is matrix scaling.

\emph{Matrix scaling.} If a nonnegative matrix $A$ has total support, there are positive
diagonal $D_1,D_2$ with $D_1AD_2$ doubly stochastic (Sinkhorn--Knopp
\cite{SinkhornKnopp}), and for symmetric $A$ one may take $D_1=D_2=D=\mathrm{diag}(d_i)$.
Then $\operatorname{per}A=\operatorname{per}(S)/\prod_id_i^2$ with $S=DAD$, and the van der
Waerden bound $\operatorname{per}S\ge n!/n^n$ \cite{Egorychev,Falikman} gives
\[
\log\operatorname{per}A\ \ge\ \log\frac{n!}{n^n}-2\sum_i\log d_i=-n-2\sum_i\log d_i+O(\log n).
\]

\emph{Bethe.} \Cref{thm:GS} with $P=S$ gives
$\log\operatorname{per}A\ge F(S)=-2\sum_i\log d_i-n+O\bigl(\sum s_{ij}^2\bigr)$. Indeed
$-\sum s_{ij}\log s_{ij}=-\sum s_{ij}(\log d_i+\log d_j)=-2\sum_i\log d_i$ because the rows and
columns sum to $1$, and $\sum(1-s)\log(1-s)=-\sum s+O(\sum s^2)=-n+O(\sum s^2)$. For
$G(n)$, $\sum s_{ij}^2\le n\max s_{ij}=O(\sqrt n)$. So van der Waerden and Bethe-at-$S$
agree up to $o(n)$. The advantage of \cref{thm:GS} is that \emph{any} doubly stochastic $P$
is allowed, so an exact scaling is not needed. This is what makes \cref{prob:L1} a
perturbation problem.

\emph{Discrete versus continuum.} Suppose $d_i=n^{-1/4}a(i/n)(1+o(1))$ in an averaged sense.
Then
\[
-2\sum_i\log d_i=\frac n2\log n-2\sum_i\log a(i/n)+o(n)=\frac n2\log n+Cn+o(n),
\]
and hence $F(S)=\frac n2\log n+(C-1)n+o(n)$. This is exactly the first-moment value of
\cref{prop:firstmoment}, since $\log n!-\frac n2\log n=\frac n2\log n-n+O(\log n)$. The
ansatz $n^{1/4}d_i\approx a(i/n)$ is consistent because the scaling equation
$d_i\sum_{j\sim i}d_j=1$ becomes, by \cref{lem:squaresum},
$a(x)\cdot\sqrt n\int a\,\kappa(x,\cdot)\cdot n^{-1/2}=1$, which is \eqref{eq:schrod}. The
same computation for the explicit ansatz \eqref{eq:ansatz} gives
$-\sum\widetilde p\log\widetilde p=\frac12\log n\sum\widetilde p-2\sum_i\log a(i/n)\sum_j\widetilde p_{ij}
=\frac n2\log n+Cn+o(n)$, since its row sums are $1+o(1)$ on average.

\emph{Where the residues went.} In the limit on $\Omega$, the analogue of
\eqref{eq:schrod} uses the kernel $\kappa(x,y)$ times the arithmetic coupling $\nu$ of
$(\alpha,\beta)$, $\beta=s-\alpha$. If $\alpha$ is Haar-distributed and $s\sim\sigma$ is
independent, then $\beta$ is Haar-distributed. So $\nu$ \emph{already has uniform
marginals}, the potentials do not depend on $\alpha$, and the optimal coupling on
$\Omega^2$ is $\mu^*\otimes\nu$. The arithmetic layer contributes $0$ to $C$. In finite
terms, this is why \cref{lem:squaresum} needs no congruence conditions: the row sums of
the ansatz involve only the sizes of the neighbours. The numerical values of $F(S)$ in
\cref{sec:deterministic}, computed for the actual graphs $G(n)$, confirm this.

\subsection{Dictionary}
\begin{center}\small
\begin{tabular}{ll}
\toprule
continuum (limit object) & finite $n$ (the graph $G(n)$)\\
\midrule
kernel $\kappa(x,y)$ on $[0,1]^2$ & $\sqrt n\,A_{G(n)}$, via \cref{lem:squaresum}\\
arithmetic coupling $\nu$ on $\Zhat^2$ & residues of $i,j$ with $i+j$ a square, \cref{prop:equid}\\
permuton $\mu\in\Pi(\lambda,\lambda)$ & doubly stochastic $P$ supported on $E(G(n))$\\
$\int\log\kappa\,d\mu-\KL(\mu\|\lambda^2)$ & $F(P)-\frac n2\log n+n$, the Bethe free energy\\
Schr\"odinger potential $a(x)$ & $n^{1/4}d_i$, the scaling vector\\
Sinkhorn on $L^\infty[0,1]$ & Sinkhorn on $A_{G(n)}$\\
$C$ & $\frac1n\bigl(\log\operatorname{per}A_{G(n)}-\frac n2\log n\bigr)+1$ (conjecturally)\\
$B=\int\log f$ (Gibbs, one marginal) & Bregman's bound via the degrees $d(i)$\\
\bottomrule
\end{tabular}
\end{center}


\bibliographystyle{amsplain}
\bibliography{square_sum}
\end{document}
